\documentclass[../../../include/open-logic-section]{subfiles}

\begin{document}

\olfileid[fr]{sth}{cardinals}{hp}
\olsection{Annexe : le principe de Hume}

Dans la \olref[cp]{sec}, nous avons présenté
le principe de Cantor :
\begin{align*}
	\card{A} = \card{B}
	& \text{ ssi } A \approx B.
\intertext{Il ressemble beaucoup à ce que l'on appelle
aujourd'hui le \emph{principe de Hume}, selon lequel :}
	\fregenum{x} {F(x)} = \fregenum{x}{G(x)}
	& \text{ ssi } F \sim G
\end{align*}
Ici, «~$F \sim G$~» abrège l'affirmation
qu'il y a exactement autant de $F$ que de
$G$, autrement dit qu'on peut mettre les
$F$ en bijection avec les $G$, c'est-à-dire :
\begin{align*}
	\exists R(&\forall v\forall y(Rvy
	\lif (Fv \land Gy)) \land {}\\
	&\forall v(Fv \lif \lexists![y][Rvy])
	\land {}\\
	&\forall y(Gy \lif \lexists![v][Rvy]))
\end{align*}
Il existe toutefois une différence de type entre
le principe de Hume et celui de Cantor. Dans
l'énoncé du principe de Cantor, les variables
«~$A$~» et «~$B$~» sont des termes du premier
ordre qui désignent des \emph{ensembles}.
Dans celui du principe de Hume, «~$F$~»,
«~$G$~» et «~$R$~» ne sont \emph{pas} des termes
du premier ordre : ils occupent plutôt une
\emph{position de prédicat}. (Peut-être désignent-ils
des \emph{propriétés}.) On pourrait donc
paraphraser le principe de Hume ainsi :
le nombre des $F$ est égal au nombre des $G$
si et seulement si les $F$ et les $G$ sont
en bijection. On l'appelle \emph{principe de Hume}
parce que Hume écrivait :
\begin{quote}
	Lorsque deux nombres sont associés de telle
	sorte qu'à chaque unité de l'un corresponde
	toujours une unité de l'autre, nous les
	déclarons égaux.
	\citep[Pt.III Bk.1 \S1]{Hume1740}
\end{quote}
Le principe de Hume a ensuite acquis une place
importante dans la logique mathématique contemporaine
grâce à \citet[\S63]{Frege1884}, qui citait ce
passage de Hume avant d'esquisser, en substance,
ce que nous appelons le principe de Hume.

Remarquez la ressemblance structurelle entre le
principe de Hume et la loi fondamentale~V.
Nous avons formulé cette dernière dans la
\olref[story][blv]{sec} ainsi :
\[
	\fregeext{x}{F(x)} = \fregeext{x}{G(x)}
	\text{ssi } \lforall[x][(F(x) \liff G(x))].
\]
Quelques explications et comparaisons peuvent
être utiles ici.

Il y a deux façons d'interpréter un principe
comme celui de Hume ou la loi fondamentale~V :
de manière \emph{prédicative} ou
\emph{imprédicative} (rappelons
la \olref[story][predicative]{sec}). Dans une
interprétation imprédicative de la loi
fondamentale~V, pour chaque~$F$, l'objet
$\fregeext{x}{F(x)}$ appartient au domaine de
quantification utilisé pour formuler cette loi
elle-même. De même, dans une interprétation
imprédicative du principe de Hume, pour
chaque~$F$, l'objet $\fregenum{x}{F(x)}$
appartient au domaine de quantification
utilisé pour formuler ce principe.
En revanche, selon l'interprétation
\emph{prédicative}, les objets
$\fregeext{x}{F(x)}$ et
$\fregenum{x}{F(x)}$ seraient des entités
appartenant à un domaine \emph{différent}.

Interprétée imprédicativement, la loi
fondamentale~V conduit à une contradiction
par la compréhension naïve (pour les détails,
voir la \olref[story][blv]{sec}). Comme la
compréhension naïve, elle peut devenir
cohérente si on l'interprète
\emph{prédicativement}. Mais elle ne fera
probablement pas tout ce que nous attendions
d'elle.

Le principe de Hume, lui, \emph{peut}
être interprété imprédicativement de façon
cohérente. Dans cette interprétation, il est
assez puissant.

Pour le voir, considérons le prédicat
«~$x \neq x$~», que rien ne satisfait,
évidemment. Le principe de Hume fournit
alors un objet $\# x( x\neq x)$, que nous
pourrions prendre pour le nombre~$0$.
Selon l'interprétation \emph{imprédicative},
et \emph{seulement} selon elle, cet objet $0$
appartient au domaine initial de quantification.
Nous pouvons donc appliquer le principe de Hume
au prédicat «~$x = 0$~» et obtenir un objet
$\#x (x = 0)$, que nous pourrions prendre pour
le nombre~$1$. De plus, le principe de Hume
entraîne que $0 \neq 1$, car il ne peut y
avoir de bijection entre les objets qui ne
sont pas identiques à eux-mêmes et ceux qui
sont identiques à $0$ : il n'y en a aucun
des premiers, mais un des seconds. Dans
l'interprétation imprédicative, $1$ appartient
lui aussi au domaine initial de quantification.
Nous pouvons donc appliquer le principe de
Hume au prédicat «~$(x = 0 \lor x = 1)$~»
pour obtenir un objet
$\#x(x = 0 \lor x = 1)$, que nous pourrions
prendre pour le nombre~$2$. Nous pouvons
alors montrer que $0\neq 2$, que
$1 \neq 2$, et ainsi de suite.

En somme, interprété imprédicativement,
le principe de Hume entraîne l'existence
d'une \emph{infinité d'objets}. C'est
pourquoi les \emph{logicistes néo-frégéens}
y ont vu un fondement de l'arithmétique.

Frege \emph{lui-même}, toutefois, n'a pas
pris le principe de Hume comme fondement
de l'arithmétique. Il l'a plutôt démontré
à partir d'une définition explicite :
$\fregenum{x}{F(x)}$ est défini comme
l'extension du concept $F \sim \Phi$.
En termes modernes, nous pourrions tenter
d'écrire cela
$\fregenum{x}{F(x)} = \Setabs{G}{F \sim G}$ ;
mais nous retrouverions alors les problèmes
de la compréhension naïve.

\end{document}
